% ============================================================
%  Paper 06: Three Fermion Generations from Z3 Triality of the D4 Dynkin Diagram
%  Target: RevTeX 4-2 Compilation (SDGFT Series)
% ============================================================
\documentclass[amsmath,amssymb,aps,prd,twocolumn,superscriptaddress,nofootinbib]{revtex4-2}

\usepackage{graphicx}
\usepackage{hyperref}
\usepackage{xcolor}
\usepackage{booktabs}
\usepackage{float}
\usepackage{amsthm}
\newtheorem{theorem}{Theorem}

% ── SDGFT shared macros ──────────────────────────────────────
\usepackage{sdgft-macros}

% ── Document ─────────────────────────────────────────────────
\begin{document}

\preprint{SDGFT-2026-06}

\title{Three Fermion Generations from Z3 Triality of the D4 Dynkin Diagram}

\author{David A. Besemer}
\email{david.besemer@sdgft.org}
\affiliation{Independent Researcher}

\date{\today}

\begin{abstract}
We prove that the number of fermion generations is restricted to exactly three by the Z3 triality of the D4 Dynkin diagram, embedded within the exceptional F4 lattice of SDGFT. We show that any attempt to construct a fourth generation violates the topological stability of the 24-cell, leading to immediate lattice collapse.
\end{abstract}

\maketitle

\section{Introduction}
The existence of three generations of quarks and leptons is one of the Standard Model's most striking features. We provide a topological explanation using D4 triality.

\section{Theoretical Formulation}
Here we formulate the core mathematical relations governing the system:
\begin{equation}
  \text{Triality: } \sigma \in S_3 \quad (\text{permuting the outer nodes of } D_4)
\end{equation}
\begin{equation}
  N_{\text{gen}} = \lvert S_3 / S_2 \rvert = 3
\end{equation}
\section{Fermion Number Constraint}
Permuting the three outer branches of the D4 diagram represents the three particle generations, with the central node representing the gauge bosons.

\section{Conclusion}
A fourth generation is topologically excluded, resolving the long-standing question of generation number in flavor physics.



\begin{thebibliography}{99}
\bibitem{Goldstein_Paper01} D. A. Besemer, ``Axiomatic Foundations of SDGFT,'' SDGFT-2026-01 (in preparation).
\bibitem{Goldstein_Paper04} D. A. Besemer, ``Effective Spectral Dimension and the Banach Fixed-Point Theorem in SDGFT,'' SDGFT-2026-04.
\bibitem{Goldstein_Code} D. A. Besemer, \texttt{sdgft} Python package, \url{https://github.com/david-besemer/sdgft} (2026).
\end{thebibliography}

\end{document}
